\begin{notechapter}
\chapter{The EM Algorithm}\label{ch:em-algorithm}
\chaptersubtitle{The three-coin model, from first principles to implementation}

\begin{chapterintro}
Many models contain variables that are never observed: the component of a
mixture that generated a point, the hidden state of a Markov chain, the coin
that was secretly chosen. Maximum likelihood is then hard, because the
likelihood is a sum inside a logarithm. The Expectation--Maximization (EM)
algorithm turns this hard problem into a sequence of easy ones. We develop it on
the smallest possible example, three coins. You will learn
\begin{itemize}[nosep]
  \item the complete-data and observed-data likelihoods, and why the latter is
        hard to maximize;
  \item the evidence lower bound (ELBO), derived with Jensen's inequality, and
        when it is tight;
  \item the E-step and M-step, with closed-form updates for the three coins;
  \item a proof that EM never decreases the likelihood, and the view of EM as
        coordinate ascent on the ELBO;
  \item what EM does on real numbers, and why the three-coin model is a
        cautionary tale about identifiability.
\end{itemize}
\prerequisites{maximum likelihood (\cref{generative-discriminative:sec:likelihood});
the Kullback--Leibler divergence and Gibbs' inequality
(\cref{bayes-decision-theory:prop:logloss}); Bayes' rule.}
\end{chapterintro}

\section{The three-coin problem}\label{em-algorithm:sec:problem}

Three coins have unknown probabilities of heads: $\pi$ for coin~A, $p$ for coin~B
and $q$ for coin~C. One experiment consists in
\begin{enumerate}[nosep]
  \item tossing coin~A;
  \item tossing coin~B if A shows heads, coin~C otherwise;
  \item recording only the result of this second toss.
\end{enumerate}
Ten experiments give $1, 1, 0, 1, 0, 0, 1, 0, 1, 1$ ($1$ = heads). We never see
which coin produced each result. Write $y_j$ for the observed result of
experiment $j$, $z_j\in\{0,1\}$ for the hidden choice ($z_j = 1$: coin~B), and
$\theta = (\pi, p, q)$. (This example appears in Li Hang's \emph{Statistical
Learning Methods}.)

\subsection{Complete and observed likelihoods}

If $z$ were observed, the probability of an outcome would be
\[
  P(y, z\mid\theta) = \bigl[\pi\,p^y(1-p)^{1-y}\bigr]^z\,\bigl[(1-\pi)\,q^y(1-q)^{1-y}\bigr]^{1-z},
\]
where the exponents select the active factor. Since $z$ is hidden, the probability
of what we see is the marginal
\[
  P(y\mid\theta) = \sum_{z\in\{0,1\}}P(y,z\mid\theta) = \pi\,p^y(1-p)^{1-y} + (1-\pi)\,q^y(1-q)^{1-y},
\]
and, by independence of the experiments, the log-likelihood of the data is
\[
  \mathcal L(\theta) = \sum_{j=1}^n\log\Bigl[\pi\,p^{y_j}(1-p)^{1-y_j} + (1-\pi)\,q^{y_j}(1-q)^{1-y_j}\Bigr].
\]

\begin{intuition}[A sum inside a logarithm]
The complete-data log-likelihood
$\log P(y,z\mid\theta) = z\log(\pi p^y(1-p)^{1-y}) + (1-z)\log((1-\pi)q^y(1-q)^{1-y})$
splits into separate terms for $\pi$, $p$ and $q$, each maximized by counting. The
observed-data log-likelihood has the sum \emph{inside} the logarithm, which ties
all parameters together: its gradient equations have no closed-form solution.
EM works with the complete-data form, replacing the unknown $z_j$ by its
expectation.
\end{intuition}

\section{A lower bound on the log-likelihood}\label{em-algorithm:sec:elbo}

Consider one observation $y$, and let $Q$ be \emph{any} distribution on the hidden
variable $z$ with $Q(z)>0$. Multiplying and dividing by $Q(z)$,
\[
  P(y\mid\theta) = \sum_zQ(z)\,\frac{P(y,z\mid\theta)}{Q(z)} = \E_{z\sim Q}\Bigl[\frac{P(y,z\mid\theta)}{Q(z)}\Bigr].
\]
Since $\log$ is concave, Jensen's inequality gives $\log\E[W]\geq\E[\log W]$, hence
\begin{equation}\label{em-algorithm:eq:elbo}
  \log P(y\mid\theta)\;\geq\;\mathrm{ELBO}(\theta, Q) = \E_Q\bigl[\log P(y,z\mid\theta)\bigr] + H(Q),
\end{equation}
where $H(Q) = -\E_Q[\log Q(z)]$ is the entropy of $Q$. The right-hand side is the
\emph{evidence lower bound}: $\log P(y\mid\theta)$ is the log-\emph{evidence}, and
the bound holds for every $Q$.

\begin{proposition}[The gap is a KL divergence]\label{em-algorithm:prop:gap}
For every $Q$,
\[
  \log P(y\mid\theta) = \mathrm{ELBO}(\theta,Q) + \mathrm{KL}\bigl(Q\,\Vert\,P(\cdot\mid y,\theta)\bigr).
\]
Hence the bound is tight if and only if $Q$ is the posterior $P(z\mid y,\theta)$.
\end{proposition}

\begin{proof}
Write $P(y,z\mid\theta) = P(y\mid\theta)\,P(z\mid y,\theta)$. Then
\[
  \mathrm{ELBO}(\theta,Q) = \E_Q\Bigl[\log\frac{P(y\mid\theta)P(z\mid y,\theta)}{Q(z)}\Bigr]
  = \log P(y\mid\theta) - \E_Q\Bigl[\log\frac{Q(z)}{P(z\mid y,\theta)}\Bigr],
\]
and the last expectation is the KL divergence. By Gibbs' inequality it is
non-negative and vanishes iff $Q = P(\cdot\mid y,\theta)$.
\end{proof}

For $n$ independent observations the bound is applied term by term, with one
distribution $Q_j$ per hidden variable $z_j$.

\section{The EM algorithm}\label{em-algorithm:sec:algorithm}

EM alternates between making the bound tight at the current parameters and
maximizing it.

\begin{center}
\begin{minipage}{0.92\linewidth}
\textbf{EM.} Start from $\theta^{(0)}$. For $t = 0, 1, 2,\dots$ until convergence:
\begin{description}[nosep]
  \item[E-step:] compute the posteriors $Q_j(z_j) = P(z_j\mid y_j,\theta^{(t)})$ and the
        expected complete-data log-likelihood
        \[
          Q(\theta,\theta^{(t)}) = \sum_{j=1}^n\E_{z_j\sim P(\cdot\mid y_j,\theta^{(t)})}\bigl[\log P(y_j,z_j\mid\theta)\bigr];
        \]
  \item[M-step:] $\theta^{(t+1)} = \argmax_\theta Q(\theta,\theta^{(t)})$.
\end{description}
\end{minipage}
\end{center}

(The entropy term of the ELBO does not depend on $\theta$, so maximizing the ELBO
over $\theta$ is maximizing $Q(\theta,\theta^{(t)})$.)

\subsection{E-step for the three coins}

The posterior that coin~B was used in experiment $j$ --- the \emph{responsibility}
of coin~B --- is given by Bayes' rule:
\[
  \gamma_j^{(t)} = P(z_j = 1\mid y_j,\theta^{(t)})
  = \frac{\pi^{(t)}(p^{(t)})^{y_j}(1-p^{(t)})^{1-y_j}}
         {\pi^{(t)}(p^{(t)})^{y_j}(1-p^{(t)})^{1-y_j} + (1-\pi^{(t)})(q^{(t)})^{y_j}(1-q^{(t)})^{1-y_j}} .
\]
Since the complete-data log-likelihood is linear in $z_j$, its expectation simply
replaces $z_j$ by $\gamma_j^{(t)}$:
\[
  Q(\theta,\theta^{(t)}) = \sum_{j=1}^n\Bigl[\gamma_j^{(t)}\log\bigl(\pi p^{y_j}(1-p)^{1-y_j}\bigr)
  + (1-\gamma_j^{(t)})\log\bigl((1-\pi)q^{y_j}(1-q)^{1-y_j}\bigr)\Bigr].
\]

\subsection{M-step for the three coins}

$Q$ separates into a term in $\pi$, a term in $p$ and a term in $q$. Setting each
derivative to zero,
\begin{align*}
  \frac{\partial Q}{\partial\pi} &= \sum_j\Bigl[\frac{\gamma_j}{\pi} - \frac{1-\gamma_j}{1-\pi}\Bigr] = 0
  &&\Longrightarrow\quad \pi^{(t+1)} = \frac1n\sum_j\gamma_j^{(t)},\\
  \frac{\partial Q}{\partial p} &= \sum_j\gamma_j\Bigl[\frac{y_j}{p} - \frac{1-y_j}{1-p}\Bigr] = 0
  &&\Longrightarrow\quad p^{(t+1)} = \frac{\sum_j\gamma_j^{(t)}y_j}{\sum_j\gamma_j^{(t)}},\\
  \frac{\partial Q}{\partial q} &= \sum_j(1-\gamma_j)\Bigl[\frac{y_j}{q} - \frac{1-y_j}{1-q}\Bigr] = 0
  &&\Longrightarrow\quad q^{(t+1)} = \frac{\sum_j(1-\gamma_j^{(t)})y_j}{\sum_j(1-\gamma_j^{(t)})} .
\end{align*}
Each update is a count with \emph{soft} assignments: $\pi$ is the average
responsibility of coin~B, and $p$ (resp.\ $q$) is the proportion of heads among the
results attributed to coin~B (resp.\ C), each result weighted by its
responsibility. Each objective is concave, so these stationary points are maxima.

\section{Why EM never decreases the likelihood}\label{em-algorithm:sec:monotone}

\begin{theorem}[Monotonicity of EM]\label{em-algorithm:thm:monotone}
$\mathcal L(\theta^{(t+1)})\geq\mathcal L(\theta^{(t)})$ for every $t$.
\end{theorem}

\begin{proof}
Let $Q^{(t)}$ denote the posteriors at $\theta^{(t)}$. By
\cref{em-algorithm:prop:gap} the bound is tight there:
$\mathcal L(\theta^{(t)}) = \mathrm{ELBO}(\theta^{(t)}, Q^{(t)})$. The M-step maximizes
$\mathrm{ELBO}(\cdot, Q^{(t)})$, and the ELBO is a lower bound for every $\theta$, so
\[
  \mathcal L(\theta^{(t+1)})\;\geq\;\mathrm{ELBO}(\theta^{(t+1)}, Q^{(t)})
  \;\geq\;\mathrm{ELBO}(\theta^{(t)}, Q^{(t)}) = \mathcal L(\theta^{(t)}). \qedhere
\]
\end{proof}

The same argument, written with $\log P(y\mid\theta) = \log P(y,z\mid\theta) - \log P(z\mid y,\theta)$
averaged over $Q^{(t)}$, shows that the gain in log-likelihood equals the gain in
$Q(\cdot,\theta^{(t)})$ \emph{plus} $\mathrm{KL}(Q^{(t)}\Vert P(\cdot\mid y,\theta^{(t+1)}))\geq0$:
the likelihood improves at least as much as $Q$ does. It is therefore enough to
\emph{increase} $Q$ in the M-step rather than maximize it (``generalized EM'').

\begin{keyidea}[EM is coordinate ascent on the ELBO]
\begin{itemize}[nosep]
  \item E-step: maximize $\mathrm{ELBO}(\theta^{(t)}, Q)$ over $Q$; by
        \cref{em-algorithm:prop:gap} the maximizer is the posterior.
  \item M-step: maximize $\mathrm{ELBO}(\theta, Q^{(t)})$ over $\theta$.
\end{itemize}
Each step can only increase the ELBO, and after each E-step the ELBO equals the
log-likelihood. When the exact posterior is intractable, restricting $Q$ to a
simpler family gives \emph{variational} EM; amortizing $Q$ with a neural network
gives the variational autoencoder. Mixtures of Gaussians and hidden Markov models
(the Baum--Welch algorithm) are classical exact instances.
\end{keyidea}

\section{EM on the data, and a lesson in identifiability}\label{em-algorithm:sec:numbers}

\Cref{em-algorithm:tab:run} shows two runs of EM on the ten results. In both, the
iteration stops changing after a single step.

\begin{table}[ht]
\centering
\caption{EM on $1,1,0,1,0,0,1,0,1,1$ from two initializations. $\gamma(1)$ and
$\gamma(0)$ are the responsibilities of coin~B for a head and a tail;
$s = \pi p + (1-\pi)q$ is the implied probability of heads.}
\label{em-algorithm:tab:run}
\begin{tabular}{@{}cccccccc@{}}
\toprule
$t$ & $\pi$ & $p$ & $q$ & $\gamma(1)$ & $\gamma(0)$ & $s$ & $\mathcal L$ \\
\midrule
0 & 0.5000 & 0.5000 & 0.5000 & 0.5000 & 0.5000 & 0.50 & $-6.931$ \\
1 & 0.5000 & 0.6000 & 0.6000 & 0.5000 & 0.5000 & 0.60 & $-6.730$ \\
\midrule
0 & 0.4000 & 0.6000 & 0.7000 & 0.3636 & 0.4706 & 0.66 & $-6.808$ \\
1 & 0.4064 & 0.5368 & 0.6432 & 0.3636 & 0.4706 & 0.60 & $-6.730$ \\
\bottomrule
\end{tabular}
\end{table}

The two runs end at different parameters with the \emph{same} likelihood. The
reason is structural. Each observation is a single toss, so its distribution is
Bernoulli with parameter
\[
  s(\theta) = P(y = 1\mid\theta) = \pi p + (1-\pi)q,
\]
and the likelihood $\mathcal L(\theta) = n_1\log s(\theta) + n_0\log(1 - s(\theta))$
depends on the three parameters only through $s(\theta)$. It is maximized by
\emph{every} $\theta$ with $s(\theta) = n_1/n = 0.6$: a two-dimensional surface of
maximum-likelihood estimates. Moreover one M-step always lands on it, because
$\pi^{(1)}p^{(1)} + (1-\pi^{(1)})q^{(1)} = \frac1n\sum_j\gamma_jy_j + \frac1n\sum_j(1-\gamma_j)y_j = \bar y$.

\begin{pitfall}[Non-identifiability is not a small-sample problem]
The parameters of the three-coin model cannot be recovered from these data,
however many experiments are performed: only $s$ is identifiable. Swapping the
roles of B and C ($(\pi,p,q)\mapsto(1-\pi,q,p)$) is a further symmetry that remains
even in identifiable models (``label switching''). Identifiability would require
more information per experiment, for instance several tosses of the selected
coin.
\end{pitfall}

\paragraph{Local maxima and restarts.} In identifiable models such as Gaussian
mixtures, EM converges (under mild conditions) to a \emph{stationary point} of the
likelihood --- usually a local maximum, sometimes a saddle point --- that depends on
the initialization. Running EM from several random starting points and keeping
the solution with the highest likelihood is standard practice.

\subsection{Implementation}

The two steps translate directly into vectorized code. The repository contains a
more complete version (\texttt{code/em-three-coin/em\_eval.py}) with random
restarts and simulated data.

\begin{lstlisting}[language=Python, caption={EM for the three-coin model.}]
import numpy as np

def em_three_coins(y, pi, p, q, n_iter=100, tol=1e-10):
    y = np.asarray(y, dtype=float)
    for _ in range(n_iter):
        # E-step: responsibility of coin B for each result
        b = pi * p**y * (1 - p)**(1 - y)
        c = (1 - pi) * q**y * (1 - q)**(1 - y)
        gamma = b / (b + c)
        # M-step: weighted counts
        new = (gamma.mean(),
               (gamma * y).sum() / gamma.sum(),
               ((1 - gamma) * y).sum() / (1 - gamma).sum())
        if max(abs(u - v) for u, v in zip(new, (pi, p, q))) < tol:
            break
        pi, p, q = new
    loglik = np.log(b + c).sum()
    return (pi, p, q), loglik

print(em_three_coins([1, 1, 0, 1, 0, 0, 1, 0, 1, 1], 0.4, 0.6, 0.7))
# ((0.4064..., 0.5368..., 0.6432...), -6.7301...)
\end{lstlisting}

\begin{chaptersummary}
\begin{itemize}[nosep]
  \item With hidden variables the log-likelihood is a log of a sum; the
        complete-data log-likelihood is a sum of logs and is easy to maximize.
  \item For any $Q$, $\log P(y\mid\theta) = \mathrm{ELBO}(\theta,Q) + \mathrm{KL}(Q\Vert P(z\mid y,\theta))$;
        the bound is tight exactly at the posterior.
  \item EM alternates computing posteriors (E) and maximizing the expected
        complete-data log-likelihood (M); for the three coins the M-step is a
        weighted count.
  \item EM never decreases the likelihood; it is coordinate ascent on the ELBO,
        the starting point of variational inference.
  \item In the three-coin model only $\pi p + (1-\pi)q$ is identifiable: EM reaches
        a maximum-likelihood estimate in one step, but which one depends on the
        initialization.
\end{itemize}
\end{chaptersummary}

\section*{Exercises}
\addcontentsline{toc}{section}{Exercises}

\begin{exercise}\label{em-algorithm:exr:1}
Verify the first E-step of the second run of \cref{em-algorithm:tab:run}: compute
$\gamma(1)$ and $\gamma(0)$ for $\theta^{(0)} = (0.4, 0.6, 0.7)$, then $\pi^{(1)}$, $p^{(1)}$
and $q^{(1)}$.
\solutionlink{em-algorithm:exr:1}
\end{exercise}

\begin{exercise}\label{em-algorithm:exr:2}
Show that if EM is started with $p^{(0)} = q^{(0)}$, then $p^{(t)} = q^{(t)} = \bar y$ and
$\pi^{(t)} = \pi^{(0)}$ for all $t\geq1$. Why is the coin~A parameter then never
learned?
\solutionlink{em-algorithm:exr:2}
\end{exercise}

\begin{exercise}\label{em-algorithm:exr:3}
Modify the experiment: for each toss of coin~A, the selected coin is tossed $m$
times and the number of heads $k_j\in\{0,\dots,m\}$ is recorded. Write the E-step
and M-step.
\solutionlink{em-algorithm:exr:3}
\end{exercise}
\end{notechapter}
